\documentclass[11pt,a4paper]{article}
\usepackage[T1]{fontenc}
\usepackage[utf8]{inputenc}
\usepackage[a4paper,margin=2.4cm]{geometry}
\usepackage{booktabs}
\usepackage[hidelinks]{hyperref}

\setlength{\parindent}{0pt}
\setlength{\parskip}{0.5em}
\renewcommand{\tablename}{Tabla}
\urlstyle{same}

\title{Evolución del proyecto de detección\\fotovoltaica en Cantabria}
\author{}
\date{Noviembre de 2026}

\begin{document}
\maketitle

El trabajo empezó como una prueba para localizar paneles solares en ortofotos de Cantabria. A la detección se añadieron una revisión de vuelos anteriores, el ajuste de pendientes y orientaciones con LiDAR y una estimación energética. Este documento recoge las pruebas que motivaron esos cambios y los resultados intermedios que quedaron registrados.

\section{Del detector francés a las etiquetas locales}

El primer ensayo reutilizó un YOLO entrenado con ejemplos franceses. El mAP50 pasó de 0,787 en la validación francesa a 0,257 en el test local. El mAP50--95 pasó de 0,468 a 0,149. Los valores se conservaron en la memoria anterior y la evaluación original ya no está disponible. Se refieren al detector inicial, mientras que las comparaciones de superficie posteriores se hacen por píxel. El ensayo motivó el etiquetado en Cantabria. Las diferencias de resolución, aspecto de las cubiertas y condiciones de captura eran explicaciones posibles, sin evidencia para aislar una de ellas.

Las ortofotos PNOA de 2023 sirvieron de referencia. Con LabelMe se dibujaron polígonos de paneles y se guardaron también imágenes revisadas sin paneles. El mismo conjunto de anotaciones se convirtió en máscaras para U-Net y en etiquetas de segmentación para YOLO. La selección incluyó imágenes aleatorias y ejemplos difíciles, entre ellos naves y cubiertas oscuras. Se reservó un test de 988 teselas, con 31 positivas.

Se compararon un YOLO11s local y dos versiones de U-Net. Se eligió la U-Net con codificador ResNet34 porque recuperaba más superficie que su versión anterior y su sesgo de área era menor que el de las otras alternativas. YOLO recuperaba más píxeles positivos, con más píxeles de fondo marcados como panel. La tabla \ref{tab:modelos} incorpora también los ensayos de noviembre sobre el mismo test.

\begin{table}[htbp]
\centering
\caption{Evaluaciones por píxel en las mismas 988 teselas del test local.}
\label{tab:modelos}
\begin{tabular}{lrrrr}
\toprule
Modelo & Dice & IoU & Exhaustividad & Sesgo de área \\
\midrule
YOLO11s local & 0,7292 & 0,5739 & 0,8658 & \(+37{,}45\,\%\) \\
U-Net anterior & 0,7944 & 0,6589 & 0,7193 & \(-18{,}89\,\%\) \\
U-Net ResNet34 elegida & 0,8560 & 0,7482 & 0,8237 & \(-7{,}54\,\%\) \\
DeepLabV3+ ResNet50 & 0,8391 & 0,7228 & 0,7893 & \(-11{,}86\,\%\) \\
U-Net ResNet34 afinada & 0,8360 & 0,7182 & 0,7815 & \(-13{,}05\,\%\) \\
\bottomrule
\end{tabular}
\end{table}

\section{Transferencia, cobertura y seguimiento}

\paragraph{Prueba con ortofotos de 2024.}
Se comparó la misma U-Net sobre 240 localizaciones con imágenes de 2023 a 15 cm/píxel y de 2024 a 20 cm/píxel. El Dice pasó de 0,8657 a 0,7254 y el sesgo de área, de \(-9{,}69\,\%\) a \(-34{,}22\,\%\). El ensayo cambió a la vez la campaña y el producto de imagen. Por ello, la caída describe la transferencia a ese conjunto y no permite atribuirla solo a la resolución. El barrido posterior de umbrales quedó como prueba exploratoria.

\paragraph{Ampliación a la cobertura regional.}
La primera selección de teselas se apoyaba en edificios. El índice regional pasó a incluir la superficie disponible fuera de ellos y a usar las huellas CBT23/LULC23 y el parcelario de 2023 como contexto. Así se conservaron, por ejemplo, detecciones en aparcamientos o suelo abierto. Tras recuperar el avance completo en noviembre, se publicaron 957\,468 teselas de 65 hojas y 18\,868 componentes detectados. Cada componente representa un polígono de la segmentación.

\paragraph{Reconstrucción de las campañas anteriores.}
Se recuperaron respuestas del visor y se cerró una revisión manual de 1\,486 imágenes de 2014, 2017 y 2020 distribuidas en 653 ventanas. Había 1\,417 imágenes positivas y 69 negativas. La agrupación espacial combinó las anotaciones con las detecciones de 2023 del lote de seguimiento y produjo 4\,118 candidatos. La muestra comparable de área de 2014 a 2020 incluye 2\,353 candidatos. Las observaciones de 2023 obtenidas del modelo se identifican como estimadas.

\section{Geometría LiDAR y estimación energética}

Se prepararon recortes LiDAR alrededor de las detecciones de 2023 del seguimiento. Se descargaron únicamente las hojas necesarias y se desarrollaron visores para comparar ortofoto, nube y planos ajustados. Los casos con retornos del suelo, varias vertientes o estructuras elevadas motivaron la comprobación con un faldón exterior. En la misma muestra de 100 casos, la versión anterior aceptó 13 pendientes y la revisada aceptó 21. Esa comparación mide cuántos casos pasaron los filtros. La exactitud angular requiere una referencia independiente.

El lote completo del seguimiento evaluó 8\,671 detecciones y aceptó 3\,973 inclinaciones. Se obtuvo azimut para 3\,280 de ellas. Las otras 693 eran suficientemente horizontales para aceptar la inclinación, con una orientación demasiado sensible a pequeñas variaciones de la normal. En esos casos se usó la orientación supuesta del escenario.

La primera estimación del inventario histórico usó ángulos supuestos. En el escenario central, el total de 2023 pasó de 35,90 GWh/año con esos ángulos a 32,20 GWh/año al incorporar la inclinación y a 31,46 GWh/año al añadir la orientación. Son comparaciones entre versiones del cálculo sobre el mismo inventario. Después se exportaron 3\,506 series por agrupación, compartiendo 2\,235 referencias PVGIS para el inventario y la meteorología de 2023.

\section{Últimos ensayos y entrega}

En la primera comparación, el conjunto exportado tenía 3\,034 imágenes de entrenamiento y 744 de validación. Después se incorporaron 60 ejemplos históricos sin paneles a entrenamiento y se revisó una selección de mil imágenes más. El conjunto exportado al cierre tenía 3\,238 imágenes de entrenamiento, 992 de validación y las mismas 988 de test. En Kaggle se probó DeepLabV3+ con ResNet50 y se afinó la U-Net. La última fila de la tabla \ref{tab:modelos} corresponde al afinado. Los dos ensayos obtuvieron menor Dice y mayor subestimación de área que la U-Net original en el test local. Por ello, el flujo de campaña conserva como referencia los pesos originales.

Quedaron disponibles la exportación regional, el inventario de seguimiento, el panel y las series horarias. Se añadió un comando que encadena las etapas y registra cuáles han terminado. Para las ortofotos de 2026 se preparó una entrada de prueba de detección. La prueba con sus imágenes queda supeditada a que estén disponibles.

\section*{Trazabilidad}

\begin{itemize}
  \item Los antecedentes del modelo francés proceden de la memoria previa. La selección de fuentes geográficas se recoge en \path{docs/revision_catastro_2023.md}.
  \item Las evaluaciones comparables están en los archivos \path{resumen.csv} de \path{runs/evaluacion/}. La comparación de 15 y 20 cm está en \path{runs/evaluacion/validacion_20cm/}.
  \item En \path{runs/indicadores/} quedan el inventario \path{seguimiento_historico/} con su resumen y la inferencia regional \path{unet_actual/} con su ficha de ejecución.
  \item El lote LiDAR se puede consultar en \path{runs/lidar/cantabria_20261112_v2/}. La producción y las series se conservan en \path{runs/visualizacion/}. Las versiones anteriores corresponden a \path{runs/visualizacion_sin_lidar/} y al CSV previo conservado en \path{runs/ejecuciones/orientacion_20261123_230217/}.
\end{itemize}

\end{document}
